\section{Results}
\label{sec:results}

Table~\ref{tab:main-results} combines condition outcomes with the three
component contrasts.  Condition C has the strongest aggregate profile:
mean regret \(0.250287\), no invalid selections or harmful non-abstentions,
and valid claims in all 18 cells.  Aggregates alone, however, obscure
concentration and do not establish uniform superiority.

Condition A records mean regret \(4.723122\), 38.89\% invalid selections,
16.67\% harmful non-abstentions, and 61.11\% claim validity.  B improves these
metrics to \(1.389345\), 11.11\%, 11.11\%, and 88.89\%, respectively.  D has
mean regret \(3.598322\), with 22.22\% invalid selections, 11.11\% harmful
non-abstentions, and 77.78\% claim validity.  These condition-level summaries
describe all intended units under canonical fail-closed accounting.

The mean and validity outcomes answer different questions.  Routing regret
aggregates the registered costs of each decision, whereas the categorical
rates show which scientific boundary was crossed.  C improves both in the
observed cells; D's higher mean accompanies invalid selections and harmful
continuations despite having no terminal execution failures.  Successful
transport is therefore not equivalent to a scientifically valid output.

% Evidence: canonical Phase 5B condition_metrics.json, paired_contrasts.json,
% uncertainty_results.json, and sensitivity_analysis.csv.
\begin{table*}[t]
\centering
\begin{tabular}{@{}lrrrrr@{}}
\toprule
Condition & Mean regret & Invalid & Harmful & Claim valid & Failed \\
\midrule
A & 4.723122 & 38.89\% & 16.67\% & 61.11\% & 1 \\
B & 1.389345 & 11.11\% & 11.11\% & 88.89\% & 1 \\
C & 0.250287 & 0\% & 0\% & 100\% & 0 \\
D & 3.598322 & 22.22\% & 11.11\% & 77.78\% & 0 \\
\midrule
\multicolumn{6}{@{}l}{\textit{Paired all-cell regret contrasts
(scenario-cluster interval)}}\\
\(B-A\) & \multicolumn{2}{l}{\(-3.333777\)}
  & \multicolumn{3}{l}{\([-4.111998,-2.555556]\)} \\
\(C-B\) & \multicolumn{2}{l}{\(-1.139058\)}
  & \multicolumn{3}{l}{\([-2.444444,0.166328]\)} \\
\(D-C\) & \multicolumn{2}{l}{\(+3.348035\)}
  & \multicolumn{3}{l}{\([1.807181,4.888889]\)} \\
\bottomrule
\end{tabular}
\caption{Canonical fail-closed outcomes, \(N=18\) per condition.  Lower
regret is better.  Intervals use only two scenario clusters and are
pilot-descriptive, not conventional significance claims.}
\label{tab:main-results}
\end{table*}


\paragraph{Instruction richness.}
Richer unified instructions reduce mean regret relative to concise prompting:
\(B-A=-3.333777\).  The direction is negative in both scenario families,
under the two-cluster interval \([-4.111998,-2.555556]\), and after trimming
(trimmed mean \(-2.235764\)).  Invalid selections fall from 38.89\% to
11.11\%, and claim validity rises from 61.11\% to 88.89\%.  This is the
clearest and most stable pilot result, without implying a population-level
effect.

The complete-pair mean is \(-3.735074\) over 16 pairs.  Scenario means are
\(-4.111998\) for GARCH and \(-2.555556\) for structural breaks.  Unlike
\(C-B\), its direction does not depend on retaining the largest absolute cell
difference.  The evidence therefore supports instruction enrichment as the
most stable explanatory component among the tested changes.

The contrast also separates instruction quality from a broad architectural
label.  A and B share the two-call unified structure and nominal output
capacity; the principal planned difference is that B explicitly assigns
planning, evidence assessment, review, and reporting responsibilities.  The
observed reduction therefore supports better task specification in this
benchmark rather than an inference about adding agents or calls.

\paragraph{Unified multi-call execution.}
The preregistered contrast favors C in the all-cell aggregate:
\(C-B=-1.139058\) (17-pair mean \(-1.206062\)).  The scenario-cluster interval
\([-2.444444,0.166328]\) crosses zero.  One severe structural-break cell
contributes 88.1\% of the total absolute contrast magnitude; excluding it
changes the mean to \(+0.088056\).  GARCH favors B
(\(+0.166328\)), whereas structural breaks favor C
(\(-2.444444\)).  C thus has the best observed aggregate profile, but the
pilot does not establish a broad staged-execution advantage.

This qualification matters despite C's favorable condition-level metrics.
Its zero invalid selections and 100\% claim validity describe the observed 18
cells, while the contrast asks what C adds over instruction-matched B.  The
answer is largely prevention of one high-loss failure rather than uniformly
lower regret on ordinary cells.  The prespecified five-cell diagnostic subset
has contrast zero, further limiting a general mechanism claim.

Under completed pairs, the mean remains negative because B's failed unit is
excluded alongside its matched C unit, but this does not resolve
concentration.  The leave-one-extreme-cell result and opposite scenario
directions show that the observed average does not reflect a uniform shift.
The correct descriptive statement is that C prevented a severe registered
failure while otherwise remaining close to B on many cells.

\paragraph{Specialized workflow.}
The tested specialized bundle performs worse than unified C:
\(D-C=+3.348035\), with interval \([1.807181,4.888889]\).  The direction
persists across both scenario families and trimming.  Because D also changes
prompt scoping, interfaces, context partitioning, and message flow, the result
does not identify persona identity as the cause and does not show that
multi-agent systems are generally inferior.

All 18 \(D-C\) pairs completed.  Scenario means are \(+1.807181\) for GARCH
and \(+4.888889\) for structural breaks, and the positive direction survives
removing the largest absolute cell.  This stability describes the tested
bundles more consistently than \(C-B\), but it cannot remove their
architectural confounding.

D's worse aggregate result is not explained by a larger nominal output
budget: C and D each allocate 8,192 output tokens across four calls.  It may
instead reflect any combination of persona scoping, intermediate contracts,
context partitioning, or coordination.  The benchmark does not distinguish
these explanations.  Consequently, the stable positive \(D-C\) contrast
characterizes the implementations without identifying role specialization as
the causal mechanism.

\paragraph{Boundary and failures.}
As a secondary Phase~4J baseline, the adaptive deterministic policy obtains
approximately \(0.0013\) mean regret.  It was co-developed with the benchmark
and directly encodes its diagnostic-acquisition and admissibility rules, so it
marks a formalizability boundary rather than a general architectural ranking.
All 72 Phase~5B units remain in the denominator; two provider-schema failures
receive the preregistered fail-closed loss treatment.  Full baseline results,
failure reconstruction, alternate interpretations, and sensitivities appear
in Supplement Sections F--H.

The two failed units occur once in A and once in B.  Canonical interpretation
assigns terminal failed loss and then subtracts matched oracle loss; it does
not invent a selected method.  Keeping these cases preserves \(N=18\) for
every condition.  Completed-pair and alternate direct-regret treatments are
reported separately and do not replace the canonical all-cell results.

The two failures also illustrate why execution reliability is part of the
scientific denominator.  Treating them as absent would compare different
cell sets and could reward a condition for failing on difficult cases.
Conversely, canonical fail-closed loss does not pretend to observe a decision
that was never emitted.  The supplement therefore retains the terminal record,
matched oracle loss, canonical loss construction, and alternate sensitivity
interpretation separately.

Under the alternate direct-regret interpretation, the failed units receive a
fixed penalty rather than terminal loss minus the matched oracle loss.  This
changes selected numerical contrasts but does not turn the concentrated
\(C-B\) result into a uniform scenario effect or reverse the stable direction
of \(B-A\) and \(D-C\).  We retain the registered interpretation as primary
and use the alternate calculation only to expose dependence on failure
accounting.
